\documentclass[11pt]{article}
\usepackage[margin=1in]{geometry}
\usepackage{amsmath, amssymb}
\usepackage{graphicx}
\usepackage{array}
\usepackage{tabularx}
\usepackage{booktabs}
\usepackage{enumitem}
\usepackage{hyperref}
\usepackage{caption}
\usepackage{titlesec}
\titleformat{\section}{\large\bfseries}{\thesection.}{1em}{}

\title{\textbf{Experiment 8: Effect of PCA on Regression and Classification Models}}
\author{}
\date{}

\begin{document}
\maketitle

%%%%%%%%%%%%%%%%%%%%%%%%%%%%%%%%%%%%%%%%%%%%%%%%%%%%%%%%%%%%
\section*{Objective}

To study the impact of \textbf{Principal Component Analysis (PCA)} 
on the performance of machine learning models.

Students must:

\begin{enumerate}[label=\alph*)]
    \item Train models in the original feature space (Without PCA)
    \item Train models in reduced feature space (With PCA)
    \item Compare performance using 5-fold cross-validation
\end{enumerate}

%%%%%%%%%%%%%%%%%%%%%%%%%%%%%%%%%%%%%%%%%%%%%%%%%%%%%%%%%%%%
\section*{Dataset}

Dataset: \textbf{PCA\_Full\_Lab\_Dataset\_1000\_Samples.csv}

\textbf{Features:} Academic and performance-related attributes  
\textbf{Regression Target:} Final\_Score\_Regression  
\textbf{Classification Target:} Performance\_Level\_Classification  

\textbf{Students must report:}
\begin{itemize}
    \item Number of samples and features
    \item Target distribution
    \item Preprocessing steps
\end{itemize}

%%%%%%%%%%%%%%%%%%%%%%%%%%%%%%%%%%%%%%%%%%%%%%%%%%%%%%%%%%%%
\section*{Preprocessing Steps}

\begin{enumerate}
    \item Handle missing values (if any)
    \item Encode categorical target (for classification)
    \item Standardize features (Very Important before PCA)
    \item Split data using 5-fold cross-validation
\end{enumerate}

%%%%%%%%%%%%%%%%%%%%%%%%%%%%%%%%%%%%%%%%%%%%%%%%%%%%%%%%%%%%
\section*{PCA Implementation}

\begin{enumerate}
    \item Apply PCA after standardization
    \item Choose:
        \begin{itemize}
            \item Either fixed number of components
            \item Or retain 95\% explained variance
        \end{itemize}
    \item Plot Scree Plot
    \item Record explained variance
\end{enumerate}

\begin{table}[h!]
\centering
\caption{PCA Summary}
\begin{tabular}{|c|c|c|}
\hline
Components Chosen & Explained Variance (\%) & Justification \\
\hline
 &  &  \\
\hline
\end{tabular}
\end{table}

%%%%%%%%%%%%%%%%%%%%%%%%%%%%%%%%%%%%%%%%%%%%%%%%%%%%%%%%%%%%
\section*{Regression Models}

\subsection*{1. Linear Regression}

\textbf{Evaluation Metric:}
\begin{itemize}
    \item Mean Squared Error (MSE)
    \item $R^2$ Score
\end{itemize}

\subsection*{2. Random Forest Regressor}

\textbf{Hyperparameters to Tune:}
\begin{itemize}
    \item Number of Trees
    \item Maximum Depth
\end{itemize}

%%%%%%%%%%%%%%%%%%%%%%%%%%%%%%%%%%%%%%%%%%%%%%%%%%%%%%%%%%%%
\section*{Regression Results Table}

\begin{table}[h!]
\centering
\caption{5-Fold CV Results – Regression}
\begin{tabularx}{\textwidth}{|X|X|X|X|X|}
\hline
Model & Metric & Fold Avg (No-PCA) & Fold Avg (With-PCA) & Observation \\
\hline
Linear Regression & MSE &  &  &  \\
Linear Regression & $R^2$ &  &  &  \\
Random Forest & MSE &  &  &  \\
Random Forest & $R^2$ &  &  &  \\
\hline
\end{tabularx}
\end{table}

%%%%%%%%%%%%%%%%%%%%%%%%%%%%%%%%%%%%%%%%%%%%%%%%%%%%%%%%%%%%
\section*{Classification Models}

\subsection*{1. Logistic Regression}

\textbf{Hyperparameter:}
\begin{itemize}
    \item Regularization parameter $C$
\end{itemize}

\textbf{Metrics:}
\begin{itemize}
    \item Accuracy
    \item F1-score
\end{itemize}

\subsection*{2. Support Vector Machine (SVM)}

\textbf{Hyperparameters:}
\begin{itemize}
    \item Kernel (Linear / RBF)
    \item C
    \item Gamma
\end{itemize}

%%%%%%%%%%%%%%%%%%%%%%%%%%%%%%%%%%%%%%%%%%%%%%%%%%%%%%%%%%%%
\section*{Classification Results Table}

\begin{table}[h!]
\centering
\caption{5-Fold CV Results – Classification}
\begin{tabularx}{\textwidth}{|X|X|X|X|X|}
\hline
Model & Metric & Fold Avg (No-PCA) & Fold Avg (With-PCA) & Observation \\
\hline
Logistic Regression & Accuracy &  &  &  \\
Logistic Regression & F1-score &  &  &  \\
SVM & Accuracy &  &  &  \\
SVM & F1-score &  &  &  \\
\hline
\end{tabularx}
\end{table}

%%%%%%%%%%%%%%%%%%%%%%%%%%%%%%%%%%%%%%%%%%%%%%%%%%%%%%%%%%%%
\section*{Analysis Questions}

\begin{enumerate}
    \item Did PCA improve Linear Regression performance? Why?
    \item Did PCA significantly affect Random Forest?
    \item Which classification model benefited more from PCA?
    \item Did PCA reduce variance across folds?
    \item Why do linear models benefit more from PCA compared to ensemble models?
\end{enumerate}

%%%%%%%%%%%%%%%%%%%%%%%%%%%%%%%%%%%%%%%%%%%%%%%%%%%%%%%%%%%%
\section*{Expected Learning Outcome}

Students should conclude:

\begin{itemize}
    \item PCA reduces multicollinearity
    \item PCA helps linear models more
    \item Tree-based models may not benefit much
    \item PCA reduces noise
    \item Dimensionality reduction can improve generalization
\end{itemize}

%%%%%%%%%%%%%%%%%%%%%%%%%%%%%%%%%%%%%%%%%%%%%%%%%%%%%%%%%%%%
\section*{Final Conclusion (To be Written by Students)}

Students must clearly state:

\begin{itemize}
    \item When PCA is useful
    \item When PCA is not necessary
    \item Trade-off between interpretability and performance
\end{itemize}

\end{document}